\chapter{Наблюдателю}
\label{ch:observer}

Режим наблюдателя нужен, когда игру хотят смотреть люди, которые в ней не участвуют: зрители стрима, друзья, будущие игроки, присматривающиеся к столу.

\section{Как смотреть}
\label{sec:observer:watch}

Мастер создаёт комнату и присылает код или готовую ссылку. Дальше --- либо откройте ссылку, либо зайдите на \texttt{/observer} и введите код.

\screenshotOptional{user/observer.png}{вход в комнату наблюдателя}

\textbf{Регистрация не нужна.} Это принципиально: ссылку можно кинуть в чат, и человек посмотрит игру, не заводя аккаунт.

\section{Что видно наблюдателю}
\label{sec:observer:visible}

Урезанный набор --- то же, что видят игроки, и ничего сверх того:

\begin{itemize}
\item текущую локацию и сцену;
\item карту с открытыми областями;
\item персонажей;
\item объект, который мастер показывает всем;
\item музыку.
\end{itemize}

Заметки мастера, скрытые области карты и личные сообщения игрокам наблюдателю не показываются. Влиять на игру он тоже не может --- только смотреть.

\section{Если ничего не происходит}
\label{sec:observer:empty}

Пустой экран у наблюдателя означает одно из трёх, и, к сожалению, выглядят эти случаи одинаково:

\begin{enumerate}
\item мастер ещё не начал игру или не открыл ни одной сцены;
\item код введён неверно;
\item комната закрыта.
\end{enumerate}

Проверьте код по символам и уточните у мастера, идёт ли игра.

\section{Мастеру: про ссылку}
\label{sec:observer:for-master}

Код комнаты --- единственная защита. Кто получил ссылку, тот может смотреть; отозвать доступ нельзя, и срок жизни у ссылки не ограничен. Не публикуйте её там, где не хотите видеть случайных зрителей.
